% ================================================================
%  第 9 章 收获 (原书 p.187-196, 边码 p.81-p.94) —— 前半
%  源文: _work/src1/task-ch09a.md
% ================================================================
\chapter{收获}

% ---------------- p.187 ----------------
\section{(33) 笔记——或新伦理}

\origpage{81}当然,一条职业伦理的规则,只有凭一种内在的态度才有意义,那态度就是它的灵魂。它无力创造它所力求表达的那种尊重与公正的态度,最多只能在一个普遍认可这条规则的圈子里,帮助这样的态度维持下去。一旦缺少了内在态度,哪怕这条规则还挂在人们嘴边,它也已失去全部意义、全部价值。

任何注释,无论多么一丝不苟、多么细致入微,也改变不了这一点。

不久前,我昔日的朋友与伙伴中有一位好心地向我解释:如今这年头,唉,数学产出涨得没边儿(这是谁都知道的),“人们”实在不得不——愿意也好,不愿意也罢——在写出来、交上来准备发表的文章里做一番严厉的筛选,只发表其中很小的一部分。他说这话时一脸真诚的遗憾,仿佛他自己也多少是这不可避免的宿命的受害者——同时,他那神情里还带点这样的意思:他自己就属于“法国那六七个人”(唉,是不幸,可事情就是这样!),由他们决定哪些文章能发表、哪些不能。我上了年纪,话少了,便只是默默听着。这个题目上有许多话可说,但我知道那会是白费口舌。一两个月后我还得知,这位同行几年前曾拒绝推荐一篇投给《法国科学院院报》(CR)的笔记,那篇笔记的作者和题目(是我七八年前向他提议的)都让我放在心上。作者为展开这个题目花了两年工夫,这题目确实不时髦(尽管我始终觉得它同样切题)。我认为他做出了出色的工作(作为第三阶段博士论文提交)。我不是这位年轻研究者的“老板”——他碰巧天分极高(我不知道他是否还会在数学里施展这些天分,看看他所受到的接待……),他的工作是在与我毫无接触的情况下做出来的。但也确实,这个题目的来源不可能有任何疑问;这可怜人一开始就处境不妙,而且他肯定还蒙在鼓里!这位同行倒还做得有礼有节,至少还有这点,我也没指望他更多,“真诚地抱歉,可您也明白……”。两%
% ---------------- p.188 ----------------
\origpage{82}年工夫——一个干劲十足的新手研究者的两年工夫,换来的却是一篇三页的《院报》笔记——那篇笔记又能花掉多少公款呢?这里有一种一眼就能看出的荒谬:两者之间那种巨大的不相称。只要肯费心去查一查深层的动机,这种荒谬肯定就消失了。只有这位同行兼旧友才能探测他自己的动机,正如只有我才能探测我的。但不必走很远我就知道,需要照顾的并不是数学产出涨得没边儿,您懂的,也不是公款(或《院报》那位想象出来的“无名读者”的耐心)……

这同一篇预备投给《院报》的笔记,先前还有幸被送到“法国那六七个人……”中的另一位手上,那位把它退给了作者的“老板”,因为这些数学“不叫他觉得好玩”(原话如此!)。(老板虽然恶心,却很谨慎,自己地位本来就不太稳,两次都宁可咽下去,也不愿得罪人……)后来我有机会跟这位同行兼旧日学生谈起此事,才知道他当初还是费心仔细读了送来的笔记,也想过一想(那笔记想必勾起他不少回忆……),并且觉得其中有些陈述本可以用对使用者更友好的方式写出来。然而他并不屑于把宝贵的时间浪费在把意见转达给当事人上:大人物的一刻钟,对一个无名年轻研究者的两年工夫!数学确实叫他“觉得好玩”到足以让他抓住这个机会,重新接触笔记里研究的那种情形(那情形一定在他心里、也在我心里,唤起了各式各样丰富的几何联想),把给出的描述吸收进去,然后——以他的家底与本事,毫不费力地——挑出其中的笨拙与疏漏。他没有浪费时间:他对某个数学情形的知识,由于一个初出茅庐的研究者两年的认真工作,变得更精确、更丰富了;那工作,大师当然几天之内就能做完(只做大略,不做证明)。这些一到手,人就想起来自己是谁了——案子已经判了,无名先生两年的工夫,正好进垃圾桶……

有些人,这种风一吹过来,什么也感觉不到——可我今天仍然被它吹得透不过气。在那件事里,这肯定正是所追求的效果之一(瞧瞧那拒绝所披的漂亮外衣),但肯定不是唯一的效果。就在同一次谈话里,这位旧友带着一种谦逊的得意向我吐露:他只有“所陈述的结果叫他吃惊,或者他不知道该怎么证明”时,才肯把一篇笔记推荐给《院报》\fnoteP{1}{(27)}{\emph{年轻人的“势利眼”,或纯洁的捍卫者}\par Ronnie Brown 向我转述了 J.H.C.怀特黑德(J.H.C. Whitehead)的一句话(布朗是他的学生),说的是“年轻人的势利眼:他们以为,一个定理的证明若是 trivial,定理便是 trivial 的”。我昔日的许多朋友,不妨好好琢磨琢磨这话。这种“势利眼”在今天绝不限于年轻人,我认识不止一位有声望的数学家,就常常如此。我对此格外敏感,因为我在数学里(别处也一样……)做得最好的东西,那些我认为最富成果的概念与结构,以及我凭耐心而顽强的工作从中抽出的本质性质,统统落进了“trivial”这个形容词之下。(这些东西,若作者不是早已成了名人,如今哪一样也难被接受为一篇《院报》笔记!)我做数学家的毕生野心,或者不如说我的激情与快乐,一直是找到那些显而易见的东西;在本书里(包括在眼下这个导论性的章节里……)我唯一的野心也是如此。决定性的东西,往往首先在于看见那个尚未被看见的问题(无论答案如何,也无论答案是否已经找到),或者抽出一个陈述(哪怕是猜想)——它概括并包含一个尚未被看见、或未被理解的局面;至于它是否被证明,证明是否 trivial,完全无关紧要,是彻底的枝节,哪怕一个仓促的、暂定的证明后来被发现是错的。怀特黑德所说的势利眼,是那种饱食终日的浪荡子的势利眼:他非要先弄清一瓶酒贵得离谱,才肯承认它好。这些年里,不止一次,我被旧日的激情一阵阵攫住,捧出我最好的东西,却看着它被那种自命不凡打回来。我感到一种至今犹新的痛,一份快乐落了空——但我还不至于流落街头,而且幸而我并不想把自己写的一篇文章安插出去。\par 怀特黑德所说的势利眼,是一种滥用权力,一种不诚实,而不只是对事物之美的麻木或封闭——当它由一个有权势的人施加于一个任其摆布的研究者身上,那人完全有条件把后者的想法吸收过来、据为己有,同时却以“显而易见”或“trivial”、因而“没有意思”为借口,堵住它们的发表。我这里甚至没有想到通常意义上的抄袭那种极端情形,那在数学界想必还很罕见。然而从实际的角度看,对付出代价的研究者来说,情形是一样的,而使它成为可能的那种内在态度,在我看来也相去不远。只不过它更舒服些:它伴随着一种对他人无限优越的感觉,伴随着那种以数学不可触碰的纯洁之毫不妥协的捍卫者自居者的问心无愧与内心满足。} (27)。这大概正是他发表得很少的一个原因。假如他把自己的标准用在自己身上,%
% ---------------- p.189 ----------------
\origpage{83}他索性就什么也不发表了。(确实,以他眼下的处境,他根本不需要发表。)他什么都晓得,要叫他吃惊,大概就跟找到一条虽可证明、他却不知如何证明的命题一样难。(二十年里,这两样我大约只碰上过两三回,而且最近十年、十五年里一次也没有!)他显然为自己的“质量”标准自豪,这标准使他在数学家的行当里成了把要求推到极致的冠军。我在其中看到的是一种百毒不侵的自我陶醉,而且不止一次看到,在那种笑眯眯、随和厚道的外表后面,是对他人毫无节制的轻蔑。我也看得出,他从中得到了极大的满足。

这位同行的情况,是我在“新伦理”的代表人物中遇到的最极端的一例。但它仍然很典型。在这里,无论是我讲的那件事,还是使之合理化的那篇信念宣言,用常识来衡量都透着一种乌布式的荒诞——其尺度如此之大,以至于这位脑子如此出众的旧友,以及肯定还有许多地位不那么显赫的同行(他们只消不找他推荐笔记就行了),都已经看不见它了。因为要看,至少得去看。只要肯费心去看动机(首先是他自己的动机),荒诞就暴露在光天化日之下,而它们同时也不再是荒诞的——它们交出了自己那份朴素的、一目了然的含义。

如果说这些年里,面对某些态度、尤其是某些行为,我常常痛苦到这般地步,那肯定是因为我隐约在其中辨认出——像是被推到极致、直到可笑或可憎地步的漫画——一些曾经属于我自己的态度与行为,它们又通过我的某位旧学生或旧友回到我身上。不止一次,我心里那套老反射又启动了:去揭发,去同那被清清楚楚指认出来的“恶”作斗争——但即便我偶尔在这里那里顺从了它,那也是带着分裂的确信。说到底,我很清楚,去斗争,仍旧是在事物的表面打滑,仍旧是回避。我的角色不是揭发,甚至也不是“改善”我所身处其中的这个世界,或“改善”我自己这个人。我的天职是学习,是通过我自己去认识这个世界,并通过这个世界认识我自己。如果我的生命能给我自己或他人带来什么好处,那也只是在我懂得忠于这份天职、懂得与自己和谐的限度之内。是时候提醒自己这一点了,好切断我心里那些老机制——它们在这里想推着我去为某个事业辩护(就说是某种已死的伦理吧),或者去说服人(去说服人相信取代它的那种伦理大概“荒诞”吧),而不是去探测、去发现、去认识,或者把描述当作探测的手段。写前面那两三页时,我除了就今日取代昨日之种种的流行态度说几句话之外,并无更确切的目的,却始终感到自己在提防着自己,那种心境,就像一个人随时准备大笔一挥,把刚写完的东西统统划掉、扔进纸篓!不过我还是要把写下的保留下来:它并不假,可它终究造出了一种虚假的局面,因为我在其中牵连别人多于牵连自己。我心底感到,我写着的时候什么也没学到——肯定正是这一点造成了我的不适。总之,是时候回到一种更实在的反思了,一种教我自己的反思,而不是妄图去教别人或说服别人的反思\fnoteP{2}{(28)}{写前面那几页时,起先我夹在“把肚子里的东西倒空”的欲望与一种保留、一种谨慎之间。所以我一直停在差不多上,这肯定是我那份不适、那种“我什么也没学到”之感的主要原因。自从记下这份不适的那几行写出之后,我把这几页改写了两遍——它们此前让我内心不满——写得更清楚地把自己摆进去,也更往事情深处走。一路下来,我当真“学到了点什么”,而且我相信,同时我还成功地用手指触到了一件重要的东西,它既超出这一具体事例,也超出我自己这个人。} (28)。

% ---------------- p.190 ----------------
\section{(34) 淤泥与源泉}

\origpage{85}我觉得,大体上,我已经把自己同其他数学家的关系——各种年纪、各种等级的数学家,那是我还属于他们那个世界、属于数学家的世界的时候——都过了一遍;同时,而且尤其是,也把我以自身态度与行为在其中所占的那一份过了一遍:那是我在今天所看到的某种精神里所占的一份,而这精神肯定不是昨天才有的。在这番反思里,或者更确切地说,在这趟旅行里,我四次遇到一些情形,它们在我看来是我身上某些态度与暧昧的典型:在那里,对他人那种自发的善意与尊重,被一些自我性的力量搅乱了,甚至一扫而光,尤其是(至少在这四例中的三例里)被一种自得搅乱了。这自得主要依仗的是这样一种所谓优越:某种脑力,加上我在数学活动上那种不成比例的投入,据说赋予了我这种优越。它从一种普遍的共识中得到确认与支撑——那共识几乎毫无保留地推崇这种脑力和这种不成比例的投入。

在这几例情形中,最后一例——那个“踩进我的地盘的没规矩的年轻人”——我觉得对眼下的论述最为重要。前三例是我这个人的典型,或我这个人的某些侧面的典型,是在某个时期(也确实是在某种环境里)——但是,如我已经有机会一再说过的,我根本不认为它们是我当时所属那个圈子的典型。我也不认为它们是今天法国数学界的典型,这么说吧——比如说,我与“不知疲倦的朋友”之间那种关系所特有的那种慢性迷误,如今大概很少见,就像当年想必也很少见一样。相反,我在“没规矩的年轻人”这件事上的态度与行为,却是今天数学界里每天都发生的事的典型,无论往哪里看。有影响力的数学家对无名青年怀抱善意、给予尊重,如今倒成了极罕见的例外——如果这青年没福气做他的学生(即便如此也未必……),也不是一位地位相当的同事向他推荐的学生。这大概正是我1970年“觉醒”之后不久就已传到我这里的东西——那场觉醒解开了许多沉默者的舌头——但那时我听到的第一手证词对我仍然是遥远的,因为它们直接涉及的,既不是我本人,也不是我圈子里最亲近的朋友。从某个时候起(大约1976年),我开始受到不止是表面的触动:那时传到我这里的回声,或我亲眼目睹的事实,其主角是这些朋友中的某几位,甚至是一些变得显赫的旧学生;而当那些遭到恶意中伤的人是我很熟悉的人、不止一次是我的学生(1970年以后的学生,那是不用说的!),因而他们的命运牵动我的时候,就更是如此。在某些情况下,而且毫无疑问:那种缺乏善意,甚至那种炫耀式的轻蔑态度,至少因为这样一个事实而变本加厉——某位年轻研究者是我的学生,或者他冒着风险(不一定非是我的学生)去做我昔日的朋友和其他同行也乐于称之为“格罗滕迪克式的玩意儿”的事情……

70年代初,“那个没规矩的年轻人”又给我写过信,非常客气地问我(他本来完全没有义务问我任何事!)是否见怪:他想发表一个定理的证明,那定理人家告诉他是我的,却从未发表过。我记得我回信时还是像过去那样一肚子不快,大概既没说行也没说不行,还让人听出这样的意思——我没看他的证明(他本来%
% ---------------- p.191 ----------------
\origpage{86}\origpage{87}当然愿意给我看,可我全不在乎,我当时正忙着我的那些活动!):那证明对我的证明肯定毫无补益(可它至少有一点好处:白纸黑字写下来,连同陈述本身,可供数学界使用!)。这恰好说明,那场著名的“觉醒”还多么浮浅,对某些根植于自得、根植于“精英至上”态度的行为毫无影响——而同一时期,我肯定正在《生存与生活》(Survivre et Vivre)上那些情真意切的文章里、在公开辩论的发言里,揭发着这些行为……

这就很具体地回答了我先前悬置的一个问题。不如就在这里承认这条朴素的真理:这类自得的态度在我身上绝没有“一劳永逸地”被克服,我怀疑有朝一日会,除非到我死。如果说确有转变,那也不是因为虚荣消失了,而是因为(重新)出现了一种好奇——对我自己这个人、对我身上某些态度与行为的真实本性的好奇。正是凭这种好奇,我才变得多少能察觉自己身上的种种虚荣表现。这深刻地改变了某种内在的动力,也因此改变了“虚荣”的效果;所谓虚荣,就是那股常常推着我去掩盖或伪造我对现实那种健全而细腻的感知的力量,其目的是壮大我自己、把自己置于他人之上,同时却宣称恰恰相反。

也许某位读者会感到困惑,就像我有朝一日也曾困惑:在我作为数学家的生活里,虚荣那种阴险而顽强的存在(他或许有时也在自己的生活里瞥见过它),与我所谓的对数学的爱、或者说激情(它也许同样在他自己对数学、或对某个别人、别的什么事的体验里唤起回声),这两者之间有一种表面的矛盾。如果他确实困惑,那么他自身就具备所需要的一切,可以(像我从前那样)重新接触到事物本身的现实——那是他可以第一手认识的东西——而不是像一只被关在词语与概念的无尽笼子里的松鼠那样打转。

看见一汪浑水的人,难道会说水与泥是一回事?要认识那并非泥的水,只消上到源头,看一看,喝一口。要认识那并非水的泥,只消登上被太阳和风晒干的河岸,在手心里掰开、捻碎一团带沙粒的黏土。野心、虚荣,或多或少能决定一个人在生活中给某种激情——比如数学的激情——留多少地盘,能在回报丰厚时把它变得贪得无厌。但最贪得无厌的野心,本身却无力发现或认识哪怕最微小的一件事,恰恰相反!在工作的时候,当理解一点点开动、成形、加深;当混乱中一点点显出秩序,或者当原本熟悉的东西忽然露出陌生的面貌,继而令人不安,直到终于爆出一个矛盾,把一种看似不可动摇的对事物的看法掀翻——在这样的工作里,没有一丝野心或虚荣的痕迹。那时领着跳舞的,是某种来自比“自我”及其不断自我膨胀的饥饿(哪怕是“知晓”与“知识”的饥饿)远为深远之处的东西——肯定比我们这个人、甚至我们这个物种都远为深远。

那就是源泉,它在我们每个人心里。

\section{(35) 我的激情}

三大激情支配了我成年后的生活,此外还有一些性质不同的力量。我终于在这些激情中认出同一种深层冲动的三种表现;那是我心中认知冲动所走的条条道路中的三条,而在我们这个无限的世界上,呈现在它面前的路有无数条。

% ---------------- p.192 ----------------
\origpage{88}在我生命里最先出现的,是我对数学的激情。十七岁那年,刚出中学,我松开缰绳,任由一种简单的癖好去跑,它就展开成了激情,在随后的二十五年里支配了我生命的进程。我“认识”数学,远在我认识第一个女人之前(除了那个我生下来就认识的女人);如今到了成熟的年纪,我发觉它始终没有燃尽。它不再支配我的生命,我也不再声称要支配它。有时它睡去,有时我甚至以为它熄灭了,它却又悄没声息地回来,一如既往地火热。它不再像从前那样吞噬我的生命——从前我把生命交给它去吞噬。它继续在我生命里留下深深的印痕,就像他爱的女人在一个情人身上留下的印痕。

我生命里的第二种激情,是对女人的追寻。这激情常以追寻伴侣的面貌出现在我面前。直到这追寻快要结束的时候,我才懂得把两者分开:那时我知道,我所追的东西哪里也找不到,或者也知道:我就把它带在自己身上。我对女人的激情,要到我母亲去世之后才真正展开(那是我第一次恋爱五年之后,那次恋爱生了一个儿子)。就在那时,二十九岁,我成了家,又有了三个孩子。起初,对孩子的依恋是对母亲依恋不可分割的一部分,是那股从女人身上发出、把我吸引到她那里去的力量的一部分。那是这份爱的激情的果实之一。

我并没有把这两种激情在我身上的并存体验为冲突,无论起初还是后来。我必定隐隐感到了两者的深层同一,而它明确显现给我,要晚得多,是在第三种激情进入我生命之后。然而,这两种激情对我生命的影响,只能是大不相同。对数学的爱把我吸引到某个世界里,即数学对象的世界;那个世界当然有它自己的“现实”,可它不是人们的生活在其中展开的那个世界。对数学事物的切身认识,可以说没有教给我任何关于我自己的东西,更不用说关于别人的东西——朝向数学的发现之冲动,只能使我远离自己和他人。有时两三个人会在这同一种冲动中相通,但那是一种表面的相通,实际上使每个人既远离自己也远离别人。所以,对数学的激情在我生命里从来不是一种成熟的力量;我也怀疑这样的激情能促进任何人的成熟\fnoteP{3}{(29)}{我这里说的是对数学、或对另一种纯智力活动的紧张而长期的投入。反过来,这样一种激情的展开,可以是与我们身上一种被遗忘的力量重新相识的方式,是同一份不肯就范的质料较量一番的机会,也一路让我们借某种真正属于自己的东西,更新并充实我们的身份感——这样的展开,完全可能是内在旅程中、成熟过程中重要的一步。} (29)。如果说我长期以来给这激情在我的生命里留了如此不成比例的位置,那肯定也正因为:它让我得以逃避对冲突的认识,也逃避对自己的认识。

相反,性的冲动,不管我们愿不愿意,都把我们径直抛向与他人的相遇,径直抛进我们自身、也抛进对方身上那个冲突的结!而我生命里对“伴侣”的追寻,却是对无冲突的幸福的追寻——那不是认知的冲动、性的冲动,并不像我喜欢相信的那样,而是对认识他人身上和我自身之中的冲突的无尽逃避。(为了结束这虚幻的追寻,以及如影随形般伴着它的不安,这是我必须学会的两件事之一。……)幸好,人尽管逃,性总会很快把人带回来!

有一天,我不再拒不接受冲突通过我爱过的、或正爱着的女人,以及通过这些爱情所生的孩子,顽强地带给我的教训。当我终于开始%
% ---------------- p.193 ----------------
\origpage{89}\origpage{90}倾听、开始学习的时候,而且在那之后许多年里,事情一直是这样:凡我学到的东西,都是通过我爱过的或正爱着的女人学到的\fnoteP{4}{(30)}{近几年,接力的是我的孩子们,由他们来给一个有时不大情愿的学生讲授人类存在的奥秘……} (30)。直到1976年,四十八岁,对女人的追寻一直是我生命里唯一巨大的成熟之力。这成熟之所以只是在此后几年、也就是最近七年才完成,是因为我用尽手边的一切手段来保护自己、不去成熟(这手段我是从父母那里、从我认识的那些圈子里学会的)。这些手段中最有效的,就是我在数学激情上的投入。

第三种大激情进入我生命的那一天——1976年10月某个夜晚——学习的大恐惧消失了。那也是对再寻常不过的现实的恐惧,首先是对关乎我这个人的那些朴素真理的恐惧,或对关乎我所珍爱的人的真理的恐惧。说来奇怪,四十八岁之前,我从未在自己身上察觉这恐惧。我是在新激情、认知激情的新表现出现的那一夜发现它的。可以说,后者取代了那终于被认出的恐惧的位置。多年来,我在别人身上把这恐惧看得清清楚楚,却有一种奇怪的盲目,使我对自己身上的它视而不见。对看的恐惧,使我看不见这同一种对看的恐惧!我和所有人一样,牢牢依恋着某种关于自己的形象,那形象大体上从童年起就没有动过。我说的那一夜,也是这旧形象第一次塌下去的一夜。另一些与它相像的形象接替了它,凭着顽固的惯性维持几天、几个月,甚至一两年,然后在审视的目光下也塌下去。懒得去看,常常推迟这样一次新的觉醒——但看的恐惧再也没有回来。凡有好奇之处,恐惧便再无容身之地。当我对自己怀着好奇的时候,我并不更害怕自己会发现什么,就像我想知道一个数学情形的底细时一样:那时有的是一种快活的期待,有时急切,却又执拗,准备好迎接一切愿意来到面前的东西,无论预料的还是意外的——一种满怀激情的专注,守候着那些毫不含糊的迹象,凭它们在假、半真与也许的初始混乱中认出真来。

在对自己好奇之中,有爱;这爱不为任何恐惧所扰——不担心我们所看的东西不合我们所愿。说实话,对自我的爱早已在这一夜之前的几个月里静静地绽开;这一夜也是这爱取得行动形式的一夜,可以说,是变得敢作敢为的一夜,毫不客气地把行头与布景撞翻!如我所说,不久又有别的行头和布景像变戏法一样冒出来,然后轮到它们被撞翻,没有咒骂,也没有咬牙切齿……

这七年里,新激情在我生命中的种种表现,终于在我看来像是一浪接一浪起伏的波涛,像一次辽阔而安详的呼吸的一进一出。这里不是试着去描出它那条蜿蜒变化之线的地方,也不是去描出作为对位的数学激情表现之线的地方。我已放弃去安排这一种或那一种的进程——如今倒不如说是两者的这双重运动在安排我生命的进程——或者更确切地说,它就是那进程本身。

在新激情出现之前的那些月份里——孕育与丰满的月份——对女人的追寻已开始改换面貌。它开始摆脱曾浸润它的那份不安,就像一口气息终于从压迫它的重负下解脱,重新找回它自己的幅度与节律。又像一团火,因为无处出气而半明半灭地闷烧着,一阵凉风吹来,忽然舒展成噼啪作响、轻盈%
% ---------------- p.194 ----------------
\origpage{91}而活泼的火焰!

火已烧足。一种看似无法止息的饥饿得到了饱足。两三年来,那场追寻似乎确已烧尽,不留一丝灰烬,把场地让给两种激情的歌唱与对唱。其一,我少年时代的激情,三十年间一直用来让我与一个被否弃的童年分离。其二,是我成熟年纪的激情,它让我重新找回了孩子,也找回了我的童年。

\section{(36) 欲望与冥想}

我说过的那一夜,一种新激情取代了一种永远消逝的旧恐惧,也是我发现冥想的那一夜。那是我第一次“冥想”之夜,它是在一种专横而急迫的需要压迫下出现的——此前几天,我仿佛被一阵阵焦虑的浪头淹没了。也许像一切焦虑一样,那是一种“起飞”式的焦虑:它执拗地向我指出:关乎我这个人的那份朴素而显眼的现实,已经同那个维持了四十年、我却从未怀疑过的自我形象脱节。肯定有一份巨大的认知渴望,与那些可观的逃跑力量并排着,也与那种逃避焦虑、想跟从前一样安宁的欲望并排着。于是有了一场紧张的工作,持续了几个小时,直到它有了结局,而我当时还不知道正在发生的事意味着什么,更不知道我要往哪里去。在这场工作中,一个个遁词被认了出来;或者更确切地说,正是这场工作让这些遁词一个个现形,每一个都表现为一种内心确信——我终于肯费心把它白纸黑字记下来,仿佛这样才更能浸透它,而它此前一直待在一种有利的模糊里。我满心欢喜地记下它,毫无提防:它必定有本事叫我着迷——那时的心境,就是一个什么都不怀疑的人的心境,对他来说,把一条未说出口的确信白纸黑字写下来,这本身就不可辩驳地证明它真实,证明它有根有据。假如我心里没有那份不知分寸、甚至该说不知羞耻的欲望——我是说认知的欲望——我每次都会停在这个“happy end”上,而每一步确实就是在这happy end 的心境里结束的。然后,我倒霉了!我会心血来潮——天知道怎么回事、为什么——要把我刚满心满意写下的东西凑近再看一看:那不是白纸黑字写着嘛,只要重读一遍就行了!可一认真重读,傻乎乎地重读,我就觉得有点不对劲,觉得不那么清楚,咦咦!然后再费心凑近一点看,就清楚了:根本不是那么回事,可以说全是瞎编,我刚才是拿尿脬当灯笼了!每一次这局部的发现都像个大惊喜,“好家伙!这一条还真不含糊!”——一个快活的惊喜,它带着一股新的劲头,把反思重新发动起来。往前走吧,总会弄到底的,肯定马上就到,只要顺着这股劲儿接着走!小结一下,理一理……于是另一种内心确信已经升上来了,带着“事情底细”的全部外表,我们只愿意相信:这一回总该对了吧;好歹记下来,图个心安,再说记下这么中肯、这么有体会的东西本身就是乐事,不赞成的人准是心术不正,这么明显的诚意,没法做得更好了,就这样,完美!

这就是新的阶段终点,新的happy end,我本来会高高兴兴停在那里的,要不是那个调皮到极点的坏小子又闹了起来——他真是不可救药,竟又把鼻子伸进这最新的“底细”和happy end里。拦不住他,于是又开始了新的一步!

% ---------------- p.195 ----------------
\origpage{92}\origpage{93}就这样,四个小时里,阶段一个接一个过去,像一颗洋葱,我把它的层一层层剥掉(这是那一夜末尾来到我脑中的形象),终于剥到了核心——到达那条再简单不过、再显眼不过的真理;说实话,那真理本就明摆在那里,可我却好几天、好几个星期(说到底,是我一辈子)都成功地把它藏匿在这一堆互相遮掩的“洋葱层”底下。

那条朴素真理终于出现,是巨大的宽慰,是意料之外、彻彻底底的解脱。那一刻我知道,我触到了焦虑的结。这五天的焦虑确确实实被化解了、消融了,转变成刚刚在我心里成形的那种认识。焦虑不只是从我眼前消失——就像整个冥想过程中那样,也像此前五天里好几次那样;而它转变成的认识,根本不属于观念的性质,也不属于那种为了了结、为了安宁而作出的让步(就像同一夜里有几回我做过的那样);那不是一件我从外面拿过来、或取得之后添到自己身上的东西。那是十足意义上的认识,第一手的,朴素而显眼,从此成了我的一部分,就像我的血肉是我的一部分。而且,它是用清楚、毫不含糊的话表达出来的——不是一大段话,而是一句三四个字的、再傻不过的短句。这表达是刚刚进行的那场工作的最后一步;只要这最后一步没跨过去,它就一直是短暂的、可逆的。在这整场工作中,把正在成形的思想、正在出现的念头仔细地、甚至一丝不苟地表达出来,是这工作本质的一部分;而每一次重新出发,都是对刚刚走过那一步的反思——那一步是我凭刚刚写下的见证而知道的(没法把它藏进一个靠不住的记忆的迷雾里!)。

在发现与解脱那一刻之后的几分钟里,我也知道了刚刚发生的事的全部分量。我刚发现的东西,比这几天那条朴素真理还要珍贵。这东西,就是我身上的一种能力——只要我有兴趣——去认识我内心发生之事的底细,认识任何分裂、冲突处境的底细;也因此,就是完全凭自己的手段去解决我心中任何我懂得去觉察的冲突的能力。

这解决不是靠什么恩典的功效完成的——前些年我倾向于那样相信——而是靠一场紧张、执拗、细致的工作,用的是我平常的官能。如果说真有“恩典”,它不在于我们心中某个冲突突然而彻底地消失,也不在于对冲突的一种理解现成地送到我们面前(就像懒人国里烤好的鸡!),而在于这认知欲望的在场或出现\fnoteP{5}{(31)}{我这里想到的是认知欲望的“阳”的形式——那种探测、发现、为出现之物命名的欲望……正是被命名,才使出现的认识变得不可逆、不可磨灭(哪怕它后来被埋没、被遗忘,哪怕它不再起作用……)。认知欲望的“阴”的、“女性”的形式,则在于一种敞开、一种接纳,在于默默迎候一种在我们存在的更深层、思想所不能抵达之处出现的认识。这样一种敞开、以及一种一时抹去一切冲突痕迹的突然认识的出现,仍像一种恩典,它深深地打动我们,尽管其可见的效果也许短暂。不过我猜想,这种如此来到我们身上的、无字的认识,在生命某些罕见的时刻,同样不可磨灭,它的作用甚至延续到我们对它的记忆之外。} (31)。正是这欲望引导我、在几个小时里把我带到冲突的核心——正如爱的欲望让我们万无一失地找到通往所爱女人最深处的那条路。

无论是自我发现还是数学,没有欲望,一切所谓的“工作”都只是装样子,哪里也去不了。最好的情况下,它让乐在其中的人无休无止地“围着罐子转”——罐里的东西是留给饿了要吃的人的!如同世上的%
% ---------------- p.196 ----------------
\origpage{94}所有人一样,我也有欲望与饥饿都不在场的时候。当涉及对自我的认知欲望时,我对自身以及我所卷入的处境的认识便停滞不动,我行事不是凭着对事情的认识,而是听任一些根深蒂固的简单机制摆布,随之而来的种种后果也一并承担——有点像一辆由电脑、而不是由人驾驶的车。但无论是冥想还是数学,没有欲望的时候、没有这份饥饿的时候,我不会想去装模作样地“工作”。所以我不曾有过冥想哪怕几个小时,或者做数学哪怕几个小时\fnoteP{6}{(32)}{\emph{炉中百铁,或:干熬无济于事!}\par 在我还做泛函分析的时候,也就是直到1954年,我有时会没完没了地死抠一个解不出的问题,明明已经没了主意,只是在旧念头的圈子里打转,而那些念头显然已经“咬”不动了。至少有一整年就是这样,特别是关于拓扑向量空间里的“逼近问题”,那问题要到二十年后才被完全不同路数的方法解决——以我当时的地步,那些方法只能与我无缘。那时驱动我的不是欲望,而是一种固执,以及对自己内心所发生之事的无知。那是难受的一年——我一生中唯一一段做数学变得难受的时期!我需要这番经历才明白:“干熬”没有用——一旦工作到了停住的地方,一察觉停住,就得转去做别的——不妨等到更合宜的时候再回到悬置的问题上。那个时刻几乎总是不久就出现——问题自己成熟起来,我根本不必去碰它,只凭在那些看似与此毫无关系的问题上兴致勃勃地工作,就足够了。我确信,当时若一味死抠,哪怕十年我也什么也弄不出来!正是从1954年起,我在数学里养成了同时往炉子里放很多块铁的习惯。我一次只打一块,但靠着一种不断重现的奇迹,我在一块上的工作也惠及所有别的铁,它们各自等着自己的时辰。后来与冥想初次接触时,情形也一样,而且完全不是我事先有意安排的——要审视的灼热问题一天天多起来,随着反思的进行而不断增加。} (32),而没学到一点东西的;而且多半(若非说总是)是某种意外、无法预见的东西。这跟我有、而别人没有的什么官能毫无关系,只在于:我真心不想做的时候,不去装模作样地做。(正是这份“想”的力量,单独就造出了我别处谈到的那种要求:在工作中不满足于差不多,非得把一种理解——无论多么朴素——走到尽头才满意。)在要发现什么的地方,没有欲望的工作就是无意义、就是装样子,跟没有欲望地做爱一样。说实话,我从未受到过这样的诱惑:把精力浪费在假装做一件我根本不想做的事上——何况有那么多让人热血沸腾的事可做,哪怕是在该睡觉的时候睡觉(还有做梦……)。

我相信,也是在同一夜,我明白了:认知的欲望与认知、发现的能力是一回事。只要我们信任它、跟随它,就是欲望把我们一直带到我们想认识的事物的核心。也还是它让我们找到——甚至不必去找——认识这些事物最有效、也最适合我们这个人的方法。至于数学,看来自古以来写作就是一种不可或缺的手段,无论“做数学”的是谁:做数学,首先就是写\fnoteP{7}{(33)}{这并不意味着,工作中没有纸(或黑板——那不过是它的一个变体!)的那些时刻,在数学工作里就不重要。尤其在“敏感时刻”更是如此:那时一种新的直觉刚刚出现,要做的不是“零敲碎打的工作”,而是以一种更整体、更直觉的方式与它“相识”——这种相识由反思的那个非正式阶段预备。在我这里,这类思考多半在床上或散步中进行,我觉得它在整个工作时间里所占的份额相对不大。同样的观察也适用于我迄今所做的冥想工作。} (33)。在一切以智力为主的发现工作里,大概也是如此。但“冥想”——我指的是自我发现的工作——肯定不必如此。然而在我这里,直到目前,写作在冥想中一直是一种有效而不可或缺的手段。就像在数学工作中一样,它是物质载体,把
% (切片止于此句: 源文 p.196 末行 "elle est le support matériel qui fixe le", 下接 p.197)
